\documentclass[../../../include/open-logic-section]{subfiles}

\begin{document}

\olfileid[tr]{sth}{ord-arithmetic}{mult}

\olsection{Sıra Sayısı Çarpımı}

Şimdi sıra sayısı çarpımına yöneliyoruz ve buna, sıra sayısı toplamına
yaklaştığımız gibi yaklaşacağız. $\alpha$ ile~$\beta$'yı çarpmak
istediğimizi varsayalım. Bunun için genişliği $\alpha$, yüksekliği $\beta$
olan dikdörtgen bir ızgara düşünebilirsiniz. $\alpha$ ile $\beta$'nın
çarpımı, her satır boyunca ilerleyip sonra bir sonraki satıra geçerek
\ldots{} bütün ızgarayı dolaşmanın sonucudur. Başka bir deyişle, $\alpha$
ile $\beta$'nın çarpımı, $\beta$'nın \emph{her} elemanının yerine
$\alpha$'nın bir kopyasını koyarak elde edilir.

Bunu biçimselleştirmek için $\alpha$ ile $\beta$'nın Kartezyen çarpımı
üzerindeki ters sözlük sıralamasını kullanmamız yeterlidir:

\begin{defn}
Her $\alpha,\beta$ sıra sayısı için bunların çarpımı
$\alpha\ordtimes\beta=\ordtype{\alpha\times\beta,\rlexless}$'dır.
\end{defn}
\noindent
Bu tanımın iyi kurulmuş olduğunu yine doğrulamamız gerekir:

\begin{lem}\ollabel{ordtimeslessiswo}
Her $\alpha$ ve $\beta$ sıra sayısı için
$\tuple{\alpha\times\beta,\rlexless}$ bir iyi sıralamadır.
\end{lem}

\begin{proof}
Tam olarak \olref[add]{ordsumlessiswo} için olduğu gibi.
\end{proof}
\noindent
Çarpımın aşağıdaki gibi davrandığını ispatlamak da zor değildir:

\begin{lem}\ollabel{ordtimesrecursion}
Her $\alpha,\beta$ sıra sayısı için:
\begin{align*}
	\alpha \ordtimes 0 &= 0\\
	\alpha \ordtimes (\beta \ordplus 1) &= 
		(\alpha \ordtimes \beta) \ordplus \alpha\\
	\alpha  \ordtimes \beta &= 
		\supstrict_{\delta < \beta}(\alpha \ordtimes \delta) && 
		\text{$\beta$ bir limit sıra sayısı olduğunda.}
\end{align*}
\end{lem}

\begin{proof}
Alıştırma olarak bırakılmıştır.
\end{proof}

Gerçekten de toplamada olduğu gibi sıra sayısı çarpımını doğrudan bir tanım
vermek yerine bu özyineleme denklemleri aracılığıyla tanımlayabilirdik.
Yine toplamada olduğu gibi bazı davranışlar tanıdıktır:

\begin{lem}\ollabel{ordinalmultiplicationisnice}
$\alpha,\beta,\gamma$ sıra sayılarıysa:
\begin{enumerate}
	\item\ollabel{ordtimes1} $\alpha\neq 0$ ve $\beta<\gamma$ ise
	$\alpha\ordtimes\beta<\alpha\ordtimes\gamma$;
	\item\ollabel{ordtimes2} $\alpha\neq 0$ ve
	$\alpha\ordtimes\beta=\alpha\ordtimes\gamma$ ise $\beta=\gamma$;
	\item\ollabel{ordtimes3} $\alpha\ordtimes(\beta\ordtimes\gamma)
	=(\alpha\ordtimes\beta)\ordtimes\gamma$;
	\item\ollabel{ordtimes4} $\alpha\leq\beta$ ise
	$\alpha\ordtimes\gamma\leq\beta\ordtimes\gamma$;
	\item\ollabel{ordtimes5} $\alpha\ordtimes(\beta\ordplus\gamma)
	=(\alpha\ordtimes\beta)\ordplus(\alpha\ordtimes\gamma)$.
\end{enumerate}
\end{lem}

\begin{proof}
Alıştırma olarak bırakılmıştır.
\end{proof}

Dilediğiniz kadar başka sonuç ispatlayabilir (ya da kaynaklardan
bulabilirsiniz). Fakat
\olref[ord-arithmetic][add]{ordsumnotcommute} göz önüne alındığında aşağıdaki
sonuç şaşırtıcı olmamalıdır:

\begin{prop}
Sıra sayısı çarpımı değişmeli değildir:
$2\ordtimes\omega=\omega<\omega\ordtimes 2$.
\end{prop}

\begin{proof}
$2\ordtimes\omega=\supstrict_{n<\omega}(2\ordtimes n)=\omega
\in\supstrict_{n<\omega}(\omega\ordplus n)=\omega\ordplus\omega
=\omega\ordtimes 2$.
\end{proof}
\noindent
Sezgisel gerekçe yine oldukça açıktır. $2\ordtimes\omega$'yı hesaplamak için
her doğal sayının yerine iki nesne koyarsınız. Bütün ``yarım'' sayıları doğal
sayıların arasına yerleştirseniz, yani
$\Setabs{\nicefrac{n}{2}}{n\in\omega}$ üzerindeki doğal sıralamayı ele
alsanız da aynı sıra tipini elde ederdiniz. Bu biçimde söylendiğinde sıra
tipinin $\omega$'nın sıra tipiyle aynı olduğu açıktır. Buna karşılık
$\omega\ordtimes 2$'yi hesaplamak için $\omega$'nın iki kopyasını art arda
yerleştirirsiniz.

\begin{prob}
\olref[sth][ord-arithmetic][mult]{ordtimeslessiswo},
\olref[sth][ord-arithmetic][mult]{ordtimesrecursion} ve
\olref[sth][ord-arithmetic][mult]{ordinalmultiplicationisnice} sonuçlarını
ispatlayın.
\end{prob}

\end{document}
